% Chapter 02.2: ডেটা ট্রান্সমিশন মেথড ও ডেটা ট্রান্সমিশন মোড
\section{ডেটা ট্রান্সমিশন মেথড ও মোড}

\begin{learningbox}
এই টপিক শেষে শিক্ষার্থী পারবে:
\begin{itemize}
\item asynchronous, synchronous ও isochronous transmission ব্যাখ্যা করতে।
\item start bit, stop bit ও clock synchronization-এর ভূমিকা বুঝতে।
\item simplex, half duplex ও full duplex mode আলাদা করতে।
\item একটি table বা tree ব্যবহার করে transmission method compare করতে।
\item board-style comparison question-এর উত্তর লিখতে।
\end{itemize}
\end{learningbox}

\subsection{ডেটা ট্রান্সমিশন মেথড}
\subsubsection{Asynchronous Transmission}
\begin{itemize}
\item Start bit, Stop bit
\item Character by character transmission
\item ব্যবহার-বর্ণনা, সুবিধা ও অসুবিধা
\end{itemize}

\subsubsection{Synchronous Transmission}
\begin{itemize}
\item Block by block transmission
\item Clock synchronization
\item উচ্চ গতির ডেটা প্রেরণ
\end{itemize}

\subsubsection{Isochronous Transmission}
\begin{itemize}
\item Real-time audio/video transmission
\item Multimedia communication-এ ব্যবহার
\end{itemize}

\subsection{ডেটা ট্রান্সমিশন মোড}
\subsubsection{Simplex}
\begin{itemize}\item একমুখী ডেটা প্রেরণ; উদাহরণ: TV broadcasting\end{itemize}
\subsubsection{Half Duplex}
\begin{itemize}\item উভয় দিকে ডেটা যায়, কিন্তু একই সময়ে নয়; উদাহরণ: walkie-talkie\end{itemize}
\subsubsection{Full Duplex}
\begin{itemize}\item একই সময়ে উভয় দিকে data আদান-প্রদান; উদাহরণ: telephone\end{itemize}

\subsection{ডেটা ট্রান্সমিশন মোড}

\begin{definitionbox}{ডেটা ট্রান্সমিশন মোড}
ডেটা প্রেরণ ও গ্রহণের দিক এবং একই সময়ে data flow-এর মাত্রাকে transmission mode বলা হয়।
\end{definitionbox}

\subsubsection{Simplex}

Simplex mode-এ data শুধু একদিকে যায়। Receiver কোনো data ফেরত পাঠায় না।

\begin{definitionbox}{সিম্প্লেক্স}
যে transmission mode-এ data কেবল একদিকে প্রেরিত হয়, তাকে simplex mode বলে।
\end{definitionbox}

\begin{examplebox}{উদাহরণ}
TV broadcasting, radio broadcasting এবং digital notice board simplex-এর সহজ উদাহরণ।
\end{examplebox}

\subsubsection{Half Duplex}

Half duplex mode-এ দুই দিকেই data যেতে পারে, কিন্তু একই সময়ে নয়। এক পক্ষ কথা বললে অন্য পক্ষ শুনে।

\begin{definitionbox}{হাফ ডুপ্লেক্স}
যে transmission mode-এ উভয় দিকে data আদান-প্রদান সম্ভব, কিন্তু একই সময়ে নয়, তাকে half duplex বলে।
\end{definitionbox}

\begin{examplebox}{উদাহরণ}
Walkie-talkie, old-style intercom এবং কিছু shared channel communication half duplex-এর উদাহরণ।
\end{examplebox}

\subsubsection{Full Duplex}

Full duplex mode-এ একই সময়ে উভয় দিকে data আদান-প্রদান হয়। এটি সবচেয়ে কার্যকর এবং স্বাভাবিক conversation-like communication-এর জন্য উপযুক্ত।

\begin{definitionbox}{ফুল ডুপ্লেক্স}
যে transmission mode-এ উভয় দিকেই একই সময়ে data প্রেরণ ও গ্রহণ করা যায়, তাকে full duplex বলে।
\end{definitionbox}

\begin{examplebox}{উদাহরণ}
Telephone call, video chat, modern internet communication এবং two-way live conferencing full duplex-এর উদাহরণ।
\end{examplebox}

\begin{center}
\begin{tabular}{|p{0.22\textwidth}|p{0.22\textwidth}|p{0.22\textwidth}|p{0.22\textwidth}|}
\hline
\textbf{Mode} & \textbf{Direction} & \textbf{Same time?} & \textbf{উদাহরণ} \\
\hline
Simplex & একদিকে & না & TV broadcasting \\
\hline
Half duplex & দুইদিকে & না & Walkie-talkie \\
\hline
Full duplex & দুইদিকে & হ্যাঁ & Telephone \\
\hline
\end{tabular}
\end{center}

\begin{boardbox}{বোর্ড ফোকাস}
Transmission method-এর প্রশ্নে method ও mode গুলোর পার্থক্য খুব পরিষ্কার করে লিখতে হয়। Table, example এবং এক লাইনের conclusion দিলে answer strong হয়।
\end{boardbox}

\begin{recapbox}
\begin{itemize}
\item Asynchronous-এ প্রতিটি character আলাদা, synchronous-এ বড় block, isochronous-এ নির্দিষ্ট সময়ভিত্তিক data যায়।
\item Simplex একদিকে, half duplex দুইদিকে কিন্তু একই সময়ে নয়, full duplex দুইদিকে একই সময়ে।
\item Real-time media-এর জন্য isochronous খুব গুরুত্বপূর্ণ।
\item পরীক্ষায় compare table আঁকলে উত্তর দ্রুত বোঝা যায়।
\end{itemize}
\end{recapbox}

\begin{practicebox}
\textbf{MCQ}
\begin{enumerate}
\item Start bit ও stop bit কোন transmission method-এ ব্যবহৃত হয়?\\
ক. Synchronous \quad খ. Asynchronous \quad গ. Isochronous \quad ঘ. Full duplex
\item একই সময়ে উভয় দিকে data আদান-প্রদান কোন mode-এ সম্ভব?\\
ক. Simplex \quad খ. Half duplex \quad গ. Full duplex \quad ঘ. Broadcast
\end{enumerate}

\textbf{সংক্ষিপ্ত প্রশ্ন}
\begin{enumerate}
\item Asynchronous ও synchronous transmission-এর মধ্যে দুটি পার্থক্য লেখ।
\item Half duplex mode-এর একটি বাস্তব উদাহরণ দাও।
\end{enumerate}

\textbf{সৃজনশীল অনুশীলন}
একটি online class-এ শিক্ষক ভিডিও, audio ও slide একসঙ্গে পাঠাচ্ছেন; আবার শিক্ষার্থীরা প্রশ্ন পাঠাচ্ছে। উদ্দীপকের আলোকে transmission method ও mode বিশ্লেষণ করো।
\end{practicebox}


\subsection{শ্রেণিবিন্যাস ও চিত্র}

\begin{center}
\fbox{\parbox{0.92\textwidth}{
\centering
\textbf{Tree diagram placeholder: Data Transmission}\\[4pt]
Data Transmission\\
$\rightarrow$ Method: Parallel, Serial\\
$\rightarrow$ Serial Synchronization: Asynchronous, Synchronous, Isochronous\\
$\rightarrow$ Direction Mode: Simplex, Half Duplex, Full Duplex\\
$\rightarrow$ Receiver-based Mode: Unicast, Broadcast, Multicast
}}
\end{center}

\begin{definitionbox}{Parallel transmission}
যে পদ্ধতিতে একটি character বা data word-এর একাধিক bit একই সময়ে একাধিক line দিয়ে গন্তব্যে পাঠানো হয়, তাকে parallel transmission বলা হয়। এটি কাছাকাছি ডিভাইসের জন্য দ্রুত হলেও বেশি line দরকার হয়।
\end{definitionbox}

\begin{definitionbox}{Serial transmission}
যে পদ্ধতিতে bit-গুলো একটির পর একটি ধারাবাহিকভাবে একটি channel দিয়ে পাঠানো হয়, তাকে serial transmission বলা হয়। দূরবর্তী যোগাযোগে এটি বেশি ব্যবহারযোগ্য।
\end{definitionbox}

\begin{center}
\begin{tabular}{|p{0.22\textwidth}|p{0.34\textwidth}|p{0.34\textwidth}|}
\hline
\textbf{ভিত্তি} & \textbf{Parallel} & \textbf{Serial} \\
\hline
Bit flow & একসাথে একাধিক bit & একটির পর একটি bit \\
\hline
দূরত্ব & কম দূরত্বে সুবিধাজনক & বেশি দূরত্বে সুবিধাজনক \\
\hline
ব্যয় & line বেশি, ব্যয় বেশি & line কম, ব্যয় কম \\
\hline
উদাহরণ & internal bus, পুরনো printer port & USB, network link, modem communication \\
\hline
\end{tabular}
\end{center}

\begin{mistakebox}{সাধারণ ভুল}
Synchronous transmission আর full duplex একই জিনিস নয়। Synchronous হলো timing/synchronization-এর পদ্ধতি; full duplex হলো একই সময়ে দুই দিকে data flow-এর mode।
\end{mistakebox}
